\documentclass[10pt,a4paper]{amsart}
\usepackage{amsmath,amssymb,amsthm,graphicx}
\usepackage{verbatim}
\newenvironment{reference}{\comment}{\endcomment}
%\newenvironment{reference}{}{}
\newenvironment{slogan}{\comment}{\endcomment}
\newenvironment{history}{\comment}{\endcomment}

% For commutative diagrams we use Xy-pic
\usepackage[all]{xy}

% We use 2cell for 2-commutative diagrams.
\xyoption{2cell}
% Theorem environments.
%
\theoremstyle{plain}
\newtheorem{theorem}[subsection]{Theorem}
\newtheorem{proposition}[subsection]{Proposition}
\newtheorem{lemma}[subsection]{Lemma}

\theoremstyle{definition}
\newtheorem{definition}[subsection]{Definition}
\newtheorem{example}[subsection]{Example}
\newtheorem{exercise}[subsection]{Exercise}
\newtheorem{situation}[subsection]{Situation}

\theoremstyle{remark}
\newtheorem{remark}[subsection]{Remark}
\newtheorem{remarks}[subsection]{Remarks}

\numberwithin{equation}{subsection}

% Macros
%
\def\lim{\mathop{\mathrm{lim}}\nolimits}
\def\colim{\mathop{\mathrm{colim}}\nolimits}
\def\Spec{\mathop{\mathrm{Spec}}}
\def\Hom{\mathop{\mathrm{Hom}}\nolimits}
\def\Ext{\mathop{\mathrm{Ext}}\nolimits}
\def\SheafHom{\mathop{\mathcal{H}\!\mathit{om}}\nolimits}
\def\SheafExt{\mathop{\mathcal{E}\!\mathit{xt}}\nolimits}
\def\Sch{\mathit{Sch}}
\def\Mor{\mathop{\mathrm{Mor}}\nolimits}
\def\Ob{\mathop{\mathrm{Ob}}\nolimits}
\def\Sh{\mathop{\mathit{Sh}}\nolimits}
\def\NL{\mathop{N\!L}\nolimits}
\def\CH{\mathop{\mathrm{CH}}\nolimits}
\def\proetale{{pro\text{-}\acute{e}tale}}
\def\etale{{\acute{e}tale}}
\def\QCoh{\mathit{QCoh}}
\def\Ker{\mathop{\mathrm{Ker}}}
\def\Im{\mathop{\mathrm{Im}}}
\def\Coker{\mathop{\mathrm{Coker}}}
\def\Coim{\mathop{\mathrm{Coim}}}

% Boxtimes
%
\DeclareMathSymbol{\boxtimes}{\mathbin}{AMSa}{"02}

%
% Macros for moduli stacks/spaces
%
\def\QCohstack{\mathcal{QC}\!\mathit{oh}}
\def\Cohstack{\mathcal{C}\!\mathit{oh}}
\def\Spacesstack{\mathcal{S}\!\mathit{paces}}
\def\Quotfunctor{\mathrm{Quot}}
\def\Hilbfunctor{\mathrm{Hilb}}
\def\Curvesstack{\mathcal{C}\!\mathit{urves}}
\def\Polarizedstack{\mathcal{P}\!\mathit{olarized}}
\def\Complexesstack{\mathcal{C}\!\mathit{omplexes}}
% \Pic is the operator that assigns to X its picard group, usage \Pic(X)
% \Picardstack_{X/B} denotes the Picard stack of X over B
% \Picardfunctor_{X/B} denotes the Picard functor of X over B
\def\Pic{\mathop{\mathrm{Pic}}\nolimits}
\def\Picardstack{\mathcal{P}\!\mathit{ic}}
\def\Picardfunctor{\mathrm{Pic}}
\def\Deformationcategory{\mathcal{D}\!\mathit{ef}}
\usepackage{cleveref}
\crefname{originalproposition}{original proposition}{original propositions}
\Crefname{originalproposition}{Original proposition}{Original propositions}
\makeatletter
\newlabel{algebra-proposition-characterize-formally-smooth}{{031J}{source}{}{sourcepointer-1}{}}
\newlabel{algebra-proposition-characterize-formally-smooth@cref}{{[originalproposition][0][]031J}{[1][0][]source}{}{}{}}
\makeatother
\crefname{originalreference}{original reference}{original references}
\Crefname{originalreference}{Original reference}{Original references}
\makeatletter
\newlabel{lemma-completion-etale}{{0G66}{source}{}{sourcepointer-2}{}}
\newlabel{lemma-completion-etale@cref}{{[originalreference][0][]0G66}{[1][0][]source}{}{}{}}
\makeatother
\crefname{originaltheorem}{original theorem}{original theorems}
\Crefname{originaltheorem}{Original theorem}{Original theorems}
\makeatletter
\newlabel{smoothing-theorem-popescu}{{07GC}{source}{}{sourcepointer-3}{}}
\newlabel{smoothing-theorem-popescu@cref}{{[originaltheorem][0][]07GC}{[1][0][]source}{}{}{}}
\makeatother
\title{A regular local ring of characteristic p with finite degree kappa over kappa\textasciicircum{}p is excellent iff its completion map is formally etale}
\author{The Stacks Project Authors}
\date{stacks-project:examples.tex:local-snapshot-20261012; GNU FDL 1.2 or later}
\begin{document}
\maketitle
\noindent\textit{Editorial note.} All delivered mathematical text is the author's own, copied from the pinned Stacks Project chapter file examples.tex; no statement or proof step was rewritten, reordered or removed.\par\medskip
\noindent\textit{Editorial note.} The counted proof is the author's complete argument: The forward direction delegates the formal smoothness to Neron-Popescu desingularization (Smoothing Ring Maps, Theorem smoothing-theorem-popescu) - load-bearing but named; content verbatim, no rewrite.\par\medskip
\noindent\textit{Editorial note.} The label lemma-excellent-regular-local-rings names the counted statement, and the proof environment immediately following it in the source is delivered verbatim as its proof.\par\medskip
\noindent\textit{Editorial note.} The counted statement is carried by the author's own lemma environment in the chapter source.\par\medskip
\noindent\textit{Allowed standard prerequisites (editorial).}\begin{itemize}
\item Regular local rings of characteristic p, the Frobenius Twist [kappa : kappa\textasciicircum{}p], and excellence
\item Formal etaleness = formal unramifiedness plus formal smoothness; Omega\_\{A\textasciicircum{}/A\} = 0 detects formal unramifiedness
\item The Neron-Popescu desingularization (Smoothing Ring Maps, theorem-popescu): the completion of an excellent G-ring is a filtered colimit of smooth A-algebras; NL and the vanishing criterion for formal smoothness (Algebra, proposition-characterize-formally-smooth)
\end{itemize}\medskip
\ifdefined\hypertarget\hypertarget{sourcepointer-1}{}\fi
\noindent\textit{Source reference.} Original proposition 031J: sources/stacks-chow/tags.txt, source lines 3986--3986. Prerequisite pointer invoked at 3602; it names a standard fact of a sibling chapter that the counted argument uses, and every step of the counted proof built on it is retained verbatim in the delivered body.\par\medskip
\ifdefined\hypertarget\hypertarget{sourcepointer-2}{}\fi
\noindent\textit{Source reference.} Original reference 0G66: sources/stacks-examples/pinned-original-source.tex, source lines 3484--3484. Prerequisite pointer invoked at 3590; it names a standard fact of the same chapter that the counted argument uses, and every step of the counted proof built on it is retained verbatim in the delivered body.\par\medskip
\ifdefined\hypertarget\hypertarget{sourcepointer-3}{}\fi
\noindent\textit{Source reference.} Original theorem 07GC: sources/stacks-chow/tags.txt, source lines 9512--9512. Prerequisite pointer invoked at 3599; it names a standard fact of a sibling chapter that the counted argument uses, and every step of the counted proof built on it is retained verbatim in the delivered body.\par\medskip
\begin{lemma}
\label{lemma-excellent-regular-local-rings}
Let $(A, \mathfrak m, \kappa)$ be a regular local ring of characteristic
$p > 0$. Suppose $[\kappa : \kappa^p] < \infty$. Then $A$ is excellent
if and only if $A \to A^\wedge$ is formally \'etale.
\end{lemma}

\begin{proof}
The backward implication follows from Lemma \ref{lemma-completion-etale}.
For the forward implication, note that we already know from
Lemma \ref{lemma-completion-etale}
that $A \to A^\wedge$ is formally unramified or equivalently
that $\Omega_{A^\wedge/A}$ is zero.
Thus, it suffices to show that the completion map is formally smooth when
$A$ is excellent. By N\'eron-Popescu
desingularization $A \to A^\wedge$ can be  written as a filtered
colimit of smooth $A$-algebras
(Smoothing Ring Maps, Theorem \ref{smoothing-theorem-popescu}).
Hence $\NL_{A^\wedge/A}$ has vanishing cohomology in degree $-1$.
Thus $A \to A^\wedge$ is formally smooth by Algebra, Proposition
\ref{algebra-proposition-characterize-formally-smooth}.
\end{proof}

\end{document}
